\begin{proof}[Proof of \cref{obj:lem:poisson-inverse-moments}]
\begin{enumerate}
\item For \(k\in\mathbb Z_{\ge0}\), write
\[
p_\lambda(k)=\Pr(K=k)=e^{-\lambda}\frac{\lambda^k}{k!},
\qquad
g(k)=\frac1{k+1},
\qquad
r_{\mathrm{inv}}(k)=\frac1{(k+1)(k+2)},
\]
with \(0^0=1\). For any real sequence, define its forward increment by
\[
\Delta f_{\mathrm{aux}}(k)
=f_{\mathrm{aux}}(k+1)-f_{\mathrm{aux}}(k),
\qquad
f_{\mathrm{aux}}\colon\mathbb Z_{\ge0}\to\mathbb R.
\]
Then
\[
0<g(k)\le1,\qquad 0<r_{\mathrm{inv}}(k)\le1,
\qquad
\Delta g(k)
=\frac1{k+2}-\frac1{k+1}
=-r_{\mathrm{inv}}(k).
\]
Thus \(g(K)\), \(r_{\mathrm{inv}}(K)\), and \(\Delta g(K)\) are square-integrable. % lean: Causalean.Mathlib.Probability.Poisson.shifted_reciprocal_memLp_two
% lean: Causalean.Mathlib.Probability.Poisson.shifted_reciprocal_pair_memLp_two

\item We first establish the forward-increment variance inequality for a sequence \(f_{\mathrm{aux}}\) of finite support. Expanding the square and using \(\sum_{k\ge0}p_\lambda(k)=1\) gives
\[
\begin{aligned}
\operatorname{Var}(f_{\mathrm{aux}}(K))
&=\frac12\sum_{i,j\ge0}p_\lambda(i)p_\lambda(j)
  \bigl(f_{\mathrm{aux}}(j)-f_{\mathrm{aux}}(i)\bigr)^2\\
&=\sum_{0\le i<j}p_\lambda(i)p_\lambda(j)
  \bigl(f_{\mathrm{aux}}(j)-f_{\mathrm{aux}}(i)\bigr)^2.
\end{aligned}
\]
For \(0\le i<j\), telescoping and Cauchy--Schwarz yield
\[
f_{\mathrm{aux}}(j)-f_{\mathrm{aux}}(i)
=\sum_{k=i}^{j-1}\Delta f_{\mathrm{aux}}(k),
\qquad
\bigl(f_{\mathrm{aux}}(j)-f_{\mathrm{aux}}(i)\bigr)^2
\le (j-i)\sum_{k=i}^{j-1}\bigl(\Delta f_{\mathrm{aux}}(k)\bigr)^2.
\]

The Poisson mass recurrence is
\[
(k+1)p_\lambda(k+1)=\lambda p_\lambda(k)
\qquad(k\ge0).
\]
Consequently, for every \(k\ge0\),
\[
\sum_{i=0}^{k}i\,p_\lambda(i)
=\lambda\sum_{i<k}p_\lambda(i),
\qquad
\sum_{j=k+1}^{\infty}j\,p_\lambda(j)
=\lambda\sum_{j=k}^{\infty}p_\lambda(j).
\]
Here and below the sum over \(i<k\) is empty when \(k=0\). These identities imply
\[
\begin{aligned}
&\sum_{i=0}^{k}\sum_{j=k+1}^{\infty}
p_\lambda(i)p_\lambda(j)(j-i)\\
&\quad=\lambda\left[
\left(\sum_{i=0}^{k}p_\lambda(i)\right)
\left(\sum_{j=k}^{\infty}p_\lambda(j)\right)
-
\left(\sum_{i<k}p_\lambda(i)\right)
\left(\sum_{j=k+1}^{\infty}p_\lambda(j)\right)
\right]\\
&\quad=\lambda p_\lambda(k)
\left[
\sum_{i<k}p_\lambda(i)+p_\lambda(k)
+\sum_{j=k+1}^{\infty}p_\lambda(j)
\right]
=\lambda p_\lambda(k).
\end{aligned}
\]
All these identities include \(\lambda=0\). Interchanging the nonnegative path sums now gives
\[
\begin{aligned}
\operatorname{Var}(f_{\mathrm{aux}}(K))
&\le
\sum_{0\le i<j}p_\lambda(i)p_\lambda(j)(j-i)
\sum_{k=i}^{j-1}\bigl(\Delta f_{\mathrm{aux}}(k)\bigr)^2\\
&=
\sum_{k\ge0}\bigl(\Delta f_{\mathrm{aux}}(k)\bigr)^2
\sum_{i=0}^{k}\sum_{j=k+1}^{\infty}
p_\lambda(i)p_\lambda(j)(j-i)\\
&=\lambda\,\mathbb E\bigl[(\Delta f_{\mathrm{aux}}(K))^2\bigr].
\end{aligned}
\]
The last sum is finite because \(\Delta f_{\mathrm{aux}}\) has finite support. % lean: Causalean.Mathlib.Probability.PoissonAddOnePoincare.poisson_addOne_poincare_finiteSupport

\item To extend this inequality, suppose that \(f_{\mathrm{aux}}(K)\) and \(\Delta f_{\mathrm{aux}}(K)\) are square-integrable. For \(N_{\mathrm{cut}}\in\mathbb Z_{\ge0}\), define
\[
f_{\mathrm{aux},N_{\mathrm{cut}}}(k)
=
\begin{cases}
f_{\mathrm{aux}}(k),&k<N_{\mathrm{cut}},\\
0,&k\ge N_{\mathrm{cut}},
\end{cases}
\]
and the errors
\[
\begin{aligned}
e^{\mathrm{val}}_{N_{\mathrm{cut}}}(k)
&=f_{\mathrm{aux},N_{\mathrm{cut}}}(k)-f_{\mathrm{aux}}(k),\\
e^{\mathrm{inc}}_{N_{\mathrm{cut}}}(k)
&=\Delta f_{\mathrm{aux},N_{\mathrm{cut}}}(k)
  -\Delta f_{\mathrm{aux}}(k).
\end{aligned}
\]
The value error satisfies
\[
\mathbb E\bigl[(e^{\mathrm{val}}_{N_{\mathrm{cut}}}(K))^2\bigr]
=\sum_{k\ge N_{\mathrm{cut}}}p_\lambda(k)f_{\mathrm{aux}}(k)^2
\longrightarrow0.
\]
For the increment error, the identity
\[
f_{\mathrm{aux}}(k+1)
=\Delta f_{\mathrm{aux}}(k)+f_{\mathrm{aux}}(k)
\]
implies
\[
f_{\mathrm{aux}}(k+1)^2
\le2\bigl((\Delta f_{\mathrm{aux}}(k))^2+f_{\mathrm{aux}}(k)^2\bigr).
\]
Thus the shifted sequence is square-integrable. The truncation definition also gives
\[
\bigl(e^{\mathrm{inc}}_{N_{\mathrm{cut}}}(k)\bigr)^2
\le2\bigl(f_{\mathrm{aux}}(k+1)^2+f_{\mathrm{aux}}(k)^2\bigr).
\]
For each fixed \(k\), this error is zero once \(N_{\mathrm{cut}}\ge k+2\). Dominated convergence therefore yields
\[
\mathbb E\bigl[(e^{\mathrm{inc}}_{N_{\mathrm{cut}}}(K))^2\bigr]
\longrightarrow0.
\]

Cauchy--Schwarz and expansion of the squares give
\[
\left|
\mathbb E[f_{\mathrm{aux},N_{\mathrm{cut}}}(K)]
-\mathbb E[f_{\mathrm{aux}}(K)]
\right|
\le
\sqrt{\mathbb E\bigl[(e^{\mathrm{val}}_{N_{\mathrm{cut}}}(K))^2\bigr]},
\]
\[
\begin{aligned}
&\left|
\mathbb E[f_{\mathrm{aux},N_{\mathrm{cut}}}(K)^2]
-\mathbb E[f_{\mathrm{aux}}(K)^2]
\right|\\
&\quad\le
\mathbb E\bigl[(e^{\mathrm{val}}_{N_{\mathrm{cut}}}(K))^2\bigr]
+2\sqrt{
\mathbb E[f_{\mathrm{aux}}(K)^2]\,
\mathbb E\bigl[(e^{\mathrm{val}}_{N_{\mathrm{cut}}}(K))^2\bigr]},
\end{aligned}
\]
and
\[
\begin{aligned}
&\left|
\mathbb E[(\Delta f_{\mathrm{aux},N_{\mathrm{cut}}}(K))^2]
-\mathbb E[(\Delta f_{\mathrm{aux}}(K))^2]
\right|\\
&\quad\le
\mathbb E\bigl[(e^{\mathrm{inc}}_{N_{\mathrm{cut}}}(K))^2\bigr]
+2\sqrt{
\mathbb E[(\Delta f_{\mathrm{aux}}(K))^2]\,
\mathbb E\bigl[(e^{\mathrm{inc}}_{N_{\mathrm{cut}}}(K))^2\bigr]}.
\end{aligned}
\]
Hence both the variance and the expected squared increment of the truncations converge to those of \(f_{\mathrm{aux}}\). Passing to the limit in the finite-support inequality proves
\[
\operatorname{Var}(f_{\mathrm{aux}}(K))
\le\lambda\,\mathbb E[(\Delta f_{\mathrm{aux}}(K))^2].
\]
Applying this to \(g\), with the square-integrability established above, gives
\[
\operatorname{Var}(g(K))
\le\lambda\,\mathbb E[r_{\mathrm{inv}}(K)^2].
\] % lean: Causalean.Mathlib.Probability.PoissonAddOnePoincare.poisson_addOne_poincare

\item We next bound the second moment of \(g(K)\). Since \(0\le g(k)^2\le1\), for \(0\le\lambda\le1\),
\[
(1+\lambda)^2\mathbb E[g(K)^2]\le4\le8.
\]
For \(\lambda>1\), the pointwise comparison
\[
g(k)^2=\frac1{(k+1)^2}
\le\frac2{(k+1)(k+2)}
=2r_{\mathrm{inv}}(k)
\]
follows from \(k+2\le2(k+1)\). Applying the Poisson recurrence twice gives
\[
\lambda^2p_\lambda(k)r_{\mathrm{inv}}(k)=p_\lambda(k+2).
\]
The reciprocal-weighted series is summable because its terms lie between zero and \(p_\lambda(k)\). Summing the identity therefore yields
\[
\begin{aligned}
\lambda^2\mathbb E[r_{\mathrm{inv}}(K)]
&=\sum_{k\ge0}p_\lambda(k+2)\\
&=1-p_\lambda(0)-p_\lambda(1)
=1-e^{-\lambda}(1+\lambda).
\end{aligned}
\]
It follows that
\[
\lambda^2\mathbb E[g(K)^2]
\le2\lambda^2\mathbb E[r_{\mathrm{inv}}(K)]
=2\bigl(1-e^{-\lambda}(1+\lambda)\bigr)
\le2.
\]
Since \((1+\lambda)^2\le4\lambda^2\) for \(\lambda>1\), we obtain
\[
(1+\lambda)^2\mathbb E[g(K)^2]\le8.
\]
Together, the two cases establish this bound for every \(\lambda\ge0\). % lean: Causalean.Mathlib.Probability.Poisson.shifted_reciprocal_factorial_second_moment
% lean: Causalean.Mathlib.Probability.Poisson.shifted_reciprocal_sq_rate_bound
% lean: Causalean.Mathlib.Probability.Poisson.shifted_reciprocal_sq_bound

\item For the squared reciprocal pair, define
\[
r_{\mathrm{inv},4}(k)
=\frac1{(k+1)(k+2)(k+3)(k+4)}
\qquad(k\in\mathbb Z_{\ge0}).
\]
The calculation
\[
6(k+1)(k+2)-(k+3)(k+4)=5k^2+11k\ge0
\]
shows that
\[
r_{\mathrm{inv}}(k)^2\le6r_{\mathrm{inv},4}(k),
\qquad
\mathbb E[r_{\mathrm{inv}}(K)^2]
\le6\mathbb E[r_{\mathrm{inv},4}(K)].
\]
Also \(0\le r_{\mathrm{inv}}(k)^2\le1\). Thus, for \(0\le\lambda\le1\),
\[
(1+\lambda)^4\mathbb E[r_{\mathrm{inv}}(K)^2]
\le16\le96.
\]

For \(\lambda>1\), four applications of the Poisson recurrence give
\[
\lambda^4p_\lambda(k)r_{\mathrm{inv},4}(k)=p_\lambda(k+4).
\]
The reciprocal-weighted series is again dominated by the Poisson mass series, so summation gives
\[
\begin{aligned}
\lambda^4\mathbb E[r_{\mathrm{inv},4}(K)]
&=\sum_{k\ge0}p_\lambda(k+4)\\
&=1-\sum_{k=0}^{3}p_\lambda(k)\\
&=1-e^{-\lambda}
\left(1+\lambda+\frac{\lambda^2}{2}+\frac{\lambda^3}{6}\right).
\end{aligned}
\]
Consequently,
\[
\begin{aligned}
\lambda^4\mathbb E[r_{\mathrm{inv}}(K)^2]
&\le6\lambda^4\mathbb E[r_{\mathrm{inv},4}(K)]\\
&=6\left[
1-e^{-\lambda}
\left(1+\lambda+\frac{\lambda^2}{2}+\frac{\lambda^3}{6}\right)
\right]
\le6.
\end{aligned}
\]
Using \((1+\lambda)^4\le16\lambda^4\) now yields
\[
(1+\lambda)^4\mathbb E[r_{\mathrm{inv}}(K)^2]\le96.
\]
This bound therefore holds for every \(\lambda\ge0\). % lean: Causalean.Mathlib.Probability.Poisson.shifted_reciprocal_factorial_fourth_moment
% lean: Causalean.Mathlib.Probability.Poisson.shifted_reciprocal_pair_sq_bound

\item To evaluate the first moment, first take \(\lambda=0\). The Poisson masses then satisfy
\[
p_0(0)=1,\qquad p_0(k)=0\quad(k\ge1),
\]
and hence
\[
\mathbb E[g(K)]=g(0)=1.
\]
For \(\lambda>0\), the recurrence gives
\[
\lambda p_\lambda(k)g(k)=p_\lambda(k+1).
\]
Because \(0\le p_\lambda(k)g(k)\le p_\lambda(k)\), this identity can be summed:
\[
\lambda\mathbb E[g(K)]
=\sum_{k\ge0}p_\lambda(k+1)
=1-p_\lambda(0)
=1-e^{-\lambda}.
\]
Dividing by \(\lambda>0\) proves the asserted positive-rate formula. % lean: Causalean.Mathlib.Probability.Poisson.shifted_reciprocal_first_moment

\item Finally, \(1+\lambda>0\) and \(8\le2^{16}=C_0\), so
\[
\mathbb E[g(K)^2]
\le\frac8{(1+\lambda)^2}
\le\frac{C_0}{(1+\lambda)^2}.
\]
Multiplying the variance inequality by the positive quantity \((1+\lambda)^4\), and then using the reciprocal-pair bound and \(\lambda\ge0\), gives
\[
\begin{aligned}
(1+\lambda)^4\operatorname{Var}(g(K))
&\le\lambda(1+\lambda)^4\mathbb E[r_{\mathrm{inv}}(K)^2]\\
&\le96\lambda
\le C_0\lambda.
\end{aligned}
\]
Division by \((1+\lambda)^4>0\) proves
\[
\operatorname{Var}(g(K))
\le\frac{C_0\lambda}{(1+\lambda)^4},
\]
including \(\lambda=0\). % lean: CausalSmith.Stat.AnnotationRarearmFrontier.poisson_inverse_moments
\end{enumerate}
\end{proof}
